\documentclass[12pt]{exam}
\usepackage[T1]{fontenc}
\usepackage[utf8]{inputenc}
\IfFileExists{lithuanian.ldf}{\usepackage[lithuanian]{babel}}{\usepackage{babel}\babelprovide[main,import]{lithuanian}}
\usepackage{lmodern,graphicx,xcolor,tikz}
\usepackage[a4paper,left=1cm,right=1cm,top=2.5cm,bottom=1cm,footskip=12pt]{geometry}
\usepackage{amsmath,amssymb,tasks}
\newcommand{\irasolaukas}[1]{\begingroup\setlength{\fboxrule}{0.3pt}\setlength{\fboxsep}{3pt}\fcolorbox{black!35}{white}{\rule{0pt}{0.45cm}\hspace{#1}}\endgroup}
\newcommand{\sprendimas}[1]{\par\smallskip\noindent\textit{Sprendimas:}\par\nobreak\vspace{0.1cm}\noindent\begin{tikzpicture}\path[use as bounding box] (0,0) rectangle (\linewidth,{-0.7*#1});\foreach \i in {1,...,#1}{\draw[draw=black!25,line width=0.3pt] (0,{-0.7*\i}) -- (\linewidth,{-0.7*\i});}\end{tikzpicture}\par\smallskip}
\newcommand{\atsakymas}[1][9]{\par\smallskip\noindent\textit{Atsakymas:} \framebox[#1cm]{\rule{0pt}{0.5cm}}\par}
\pagestyle{headandfoot}\runningheadrule
\pointsdroppedatright\bracketedpoints\setlength{\rightpointsmargin}{1.6cm}
\pointpoints{taškas}{taškai}
\pointname{ tšk.}\marginpointname{ tšk.}
\firstpagefooter{}{\thepage}{}
\runningfooter{}{\thepage}{}
\setlength{\parindent}{0pt}
\firstpageheader{\small\begin{tabular}{@{}l@{\hspace{0.2cm}}l@{}}Mokinio kodas: & \irasolaukas{2.3cm} \\[0.15cm] Klasė: & \irasolaukas{2.3cm}\end{tabular}}{\small\textbf{Kontrolinio darbo perrašymas}\\45 min.}{\small Vilties variantas}
\runningheader{\small Kvadratinės lygtys}{}{\small Vilties variantas}
\begin{document}
\begin{center}\textbf{Kvadratinės lygtys}\end{center}
Visuose uždaviniuose laikykite, kad $x\in\mathbb R$. Sprendimus pagrįskite. Iš viso -- 40 taškų.
\medskip
\begin{questions}
\question[6] Išspręskite lygtį:\droppoints
\[(x+3)(x-2)=2x+6.\]
\sprendimas{7}\atsakymas[12]
\question[4] Lygties $2x^2-6x+1=0$ šaknys yra $x_1$ ir $x_2$. Nerasdami šaknų, apskaičiuokite $x_1^3+x_2^3$.\droppoints
\sprendimas{6}\atsakymas[12]
\newpage
\question[8] Raskite visas parametro $m$ reikšmes, su kuriomis lygties\droppoints
\[x^2-(m+2)x+2m=0\]
realiosios šaknys $x_1$ ir $x_2$ tenkina sąlygą $x_1^2+x_2^2=20$. Patikrinkite, ar gautos $m$ reikšmės užtikrina realiųjų šaknų egzistavimą.
\sprendimas{12}\atsakymas[12]
\question[7] Įvesdami tinkamą keitinį, išspręskite lygtį:\droppoints
\[x^4-5x^2-6=0.\]
\sprendimas{8}\atsakymas[12]
\newpage
\question[9] Įvesdami tinkamą keitinį, išspręskite lygtį:\droppoints
\[\frac{x-2}{x+1}-2\frac{x+1}{x-2}=1.\]
\sprendimas{9}\atsakymas[12]
\question[6] Išskaidykite dauginamaisiais:\droppoints
\begin{tasks}[label=\alph*),label-width=1.2em](1)
\task $2x^2-7x+3$;
\task $x^3-3x^2+2x-6$
\end{tasks}
\sprendimas{7}\atsakymas[12]
\end{questions}
\newpage
\begin{center}\large\textbf{Juodraštis}\end{center}
\end{document}